\chapter{Riesgos y medidas de protección}

Como ya se mencionó, la utilización de materiales radiactivos implica la presencia de alguno o varios de los riesgos siguientes:

\begin{description}
  \item[Irradiación externa:] gamma, beta.
  \item[Irradiación interna:] alfa, beta, gamma.
  \item[Contaminación externa:] gamma, beta.
  \item[Contaminación interna:] alfa, beta, gamma.
\end{description}

\section{Liberación de materiales radiactivos}

Los materiales radiactivos podrán liberarse al medio ambiente cuando ocurra algún accidente que afecte la integridad del contenedor y cápsula. Cuando esto ocurre, la presencia del material radiactivo ocasiona riesgos de contaminación externa y/o interna, ya que el material puede embarrarse en la piel durante las actividades de trabajo e incluso ingerirse o inhalarse, dependiendo de las características físicas del material radiactivo.

\section{Medidas de protección contra la irradiación externa}

A continuación se mencionan algunas medidas de protección contra la irradiación externa que se deberán seguir siempre al tomar una radiografía:

\begin{enumerate}
  \item Antes de empezar el trabajo, comprobar que el dosímetro de lectura directa no se encuentre saturado o fuera de escala y que el dosímetro de película esté en buen estado.
  \item Verificar que el detector Geiger tenga sus baterías en buen estado, que esté bien calibrado y que su alarma esté encendida.
  \item Antes de empezar la irradiación, acordonar un área aproximada de 15 m de la fuente con letreros que indiquen la presencia de radiación.
  \item Durante el tiempo que dure la exposición, mantenerse alejado de la fuente, pero controlar el área para evitar que algunas personas puedan acercarse y recibir una exposición innecesaria.
  \item Al terminar la exposición, verificar con el detector Geiger que la fuente efectivamente regresó al contenedor y que se encuentra segura. Se deben tomar los niveles de radiación cerca del contenedor y de la manguera posicionadora.
  \item Volver a leer el dosímetro de bolsillo y tomar nota de la dosis recibida durante la jornada. Al final del día, registrar la dosis en las formas previstas para ello y, al final de cada semana, entregarla al encargado.
  \item En caso de exposiciones con un grado de dificultad mayor, como por ejemplo en tuberías aéreas, en donde el técnico debe permanecer cerca de la fuente, el tiempo de irradiación será muy corto y no permitirá alejarse y regresar. El uso del colimador es indispensable, además de que la persona tendrá que planear el trabajo de forma tal que se aplique lo aprendido.
\end{enumerate}
